\BeleskeID{05-s032}
% element: 05-s032:d01
Logičko pomeranje u prazne pozicije unosi nule. Pri aritmetičkom pomeranju udesno ponavlja se bit znaka, pa drugi komplement ostaje smislen i za negativne vrednosti. Na primer, $1001$ predstavlja $-7$; logičko pomeranje daje $0100=4$, dok aritmetičko daje $1100=-4$.

Pomeranje ulevo množi vrednost sa dva samo dok tačan rezultat može da stane u registar. U drugom komplementu za jedan korak prekoračenje se prepoznaje kada se prethodna dva najviša bita razlikuju. Pri pomeranju udesno odbačeni bit može predstavljati izgubljeni razlomljeni deo: za negativne neparne brojeve aritmetički pomeraj daje količnik zaokružen prema minus beskonačnosti, a ne prema nuli.
